\documentclass{article}
\usepackage{neurips_2026}
\usepackage{enumitem}
\usepackage[utf8]{inputenc} % allow utf-8 input
\usepackage[T1]{fontenc}    % use 8-bit T1 fonts
\usepackage{hyperref}       % hyperlinks
\usepackage{url}            % simple URL typesetting
\usepackage{mathrsfs}
\usepackage{booktabs}       % professional-quality tables
\usepackage{amsfonts}       % blackboard math symbols
\usepackage{nicefrac}       % compact symbols for 1/2, etc.
\usepackage{microtype}      % microtypography
\usepackage{xcolor}         % colors
\usepackage{xcolor}


\usepackage{microtype}
\usepackage{graphicx}
\usepackage{subfigure}
\usepackage{booktabs} % for professional tables
\usepackage{bbm}
\newcommand\R{\mathbb{R}}
% \newcommand{\ishani}[1]{\textcolor{blue}{#1}}
% \newcommand{\ishanicomment}[1]{\textcolor{red}{#1}}

% hyperref makes hyperlinks in the resulting PDF.
% If your build breaks (sometimes temporarily if a hyperlink spans a page)
% please comment out the following usepackage line and replace
% \usepackage{icml2025} with \usepackage[nohyperref]{icml2025} above.
\usepackage{hyperref}

% Attempt to make hyperref and algorithmic work together better:
\usepackage{algorithm}
\usepackage[noend]{algorithmic}
\newcommand{\Comment}[1]{\textcolor{gray}{\footnotesize\textit{// #1}}}


% Use the following line for the initial blind version submitted for review:

% If accepted, instead use the following line for the camera-ready submission:
% \usepackage[accepted]{icml2025}

% For theorems and such
\usepackage{amsmath}
\usepackage{amssymb}
\usepackage{mathtools}
\usepackage{amsthm}

% if you use cleveref..
\usepackage[capitalize,noabbrev]{cleveref}

\usepackage{tikz}
\usepackage{tikz-cd}
\usetikzlibrary{automata,positioning}

%%%%%%%%%%%%%%%%%%%%%%%%%%%%%%%%
% THEOREMS
%%%%%%%%%%%%%%%%%%%%%%%%%%%%%%%%
\theoremstyle{plain}
\newtheorem{theorem}{Theorem}[section]
\newtheorem{proposition}[theorem]{Proposition}
\newtheorem{example}[theorem]{Example}

\newtheorem{lemma}[theorem]{Lemma}
\newtheorem{corollary}[theorem]{Corollary}
\theoremstyle{definition}
\newtheorem{definition}[theorem]{Definition}
\newtheorem{assumption}[theorem]{Assumption}
\newtheorem{remark}[theorem]{Remark}
\newtheorem{discussion}[theorem]{Discussion}
\newtheorem{property}[theorem]{Property}
\def\ceil#1{\left\lceil #1 \right\rceil}

\def\est{\mathrm{est}}
\def\E{\mathbb{E}}
\def\N{\mathbb{N}}
\DeclareMathOperator*{\argmax}{arg\,max}

\newcommand{\Real}{\mathbb{R}}
\newcommand{\Nat}{\mathbb{N}}
\newcommand{\F}{\mathbb{F}}
\newcommand{\Z}{\mathbb{Z}}
\renewcommand{\S}{\mathbb{S}}%

% probability
\renewcommand{\P}{\mathbb{P}}

\DeclareMathOperator*{\argmin}{argmin}
\newcommand{\poly}{\operatorname{poly}}
\newcommand{\polylog}{\operatorname{polylog}}
\newcommand{\mdiag}{\mathbf{Diag}} % diagonal matrix with vector on diag
\newcommand{\sign}{\operatorname{sign}}%
\newcommand{\nnz}{\operatorname{nnz}}%
\newcommand{\tail}{\operatorname{tail}}%
\newcommand{\dist}{\operatorname{dist}}%

% redfine some classical notation
\renewcommand{\tilde}{\widetilde}
\renewcommand{\hat}{\widehat}
\renewcommand{\bar}{\overline}

\usepackage{pgffor}
\foreach \x in {A,...,Z}{
		\expandafter\xdef\csname m\x\endcsname{\noexpand\mathbf{\x}}
	}
\foreach \x in {A,...,Z}{
		\expandafter\xdef\csname om\x\endcsname{\noexpand\overline{\noexpand\mathbf{\x}}}
	}
\foreach \x in {A,...,Z}{
		\expandafter\xdef\csname c\x\endcsname{\noexpand\mathcal{\x}}
	}


\makeatletter

\DeclareRobustCommand\widecheck[1]{{\mathpalette\@widecheck{#1}}}
\def\@widecheck#1#2{%
	\setbox\z@\hbox{\m@th$#1#2$}%
	\setbox\tw@\hbox{\m@th$#1%
			\widehat{%
				\vrule\@width\z@\@height\ht\z@
				\vrule\@height\z@\@width\wd\z@}$}%
	\dp\tw@-\ht\z@
	\@tempdima\ht\z@ \advance\@tempdima2\ht\tw@ \divide\@tempdima\thr@@
	\setbox\tw@\hbox{%
		\raise\@tempdima\hbox{\scalebox{1}[-1]{\lower\@tempdima\box
				\tw@}}}%
	{\ooalign{\box\tw@ \cr \box\z@}}}


\newcommand{\manifold}{\mathcal{M}}
\newcommand{\dataset}{\mathcal{D}}
\newcommand{\domain}{\mathcal{X}}
\newcommand{\ffam}{\mathcal{F}}
\newcommand{\relu}{\mathrm{ReLU}}
\newcommand{\vspan}{\mathrm{span}}

\newcommand{\todo}[1]{\textcolor{red}{[TODO: #1]}}
\newcommand{\nastia}[1]{\textcolor{pink}{[Nastia says: #1]}}
\newcommand{\emile}[1]{\textcolor{blue}{[Emile says: #1]}}
\newcommand{\pan}[1]{\textcolor{magenta}{[Pan says: #1]}}
\newcommand{\wenjing}[1]{\textcolor{pink}{[Wenjing says: #1]}}

\newcommand{\TV}{\mathrm{TV}} % total variation norm
\newcommand{\deff}{d_{\mathrm{eff}}} % effective dimension
\newcommand{\CD}{\mathrm{CD}} % curvature-dimension condition

\newcommand{\kernel}{\kappa}

\newcommand{\HK}{\mathcal{H}_\kernel} % Hilbert space of functions induced by a kernel $K$


% \newcommand{\HK}{\mathcal{H}_K}
\newcommand{\X}{\mathcal{X}}
\newcommand{\B}{\mathcal{B}}
% \newcommand{\E}{\mathbb{E}}
% \newcommand{\R}{\mathbb{R}}
% \newcommand{\N}{\mathbb{N}}
\newcommand{\iid}{\stackrel{\mathrm{iid}}{\sim}}
\newcommand{\inprod}[2]{\langle #1, #2 \rangle}
\newcommand{\norm}[1]{\| #1 \|}

\DeclareMathOperator{\Cov}{Cov}
\DeclareMathOperator{\Var}{Var}
\DeclareMathOperator{\Range}{Range}
\DeclareMathOperator{\Tr}{Tr}
\newcommand{\Splus}{\mathcal{S}_+^d}
\newcommand{\Splusplus}{\mathcal{S}_{++}^d}





\makeatother

% % \usepackage[bookmarks,colorlinks,breaklinks]{hyperref}
\hypersetup{urlcolor=blue, colorlinks=true, citecolor=green!50!black, linkcolor=blue}
\usepackage{multirow}
\usepackage[T1]{fontenc} % For polish hook https://tex.stackexchange.com/questions/75396/how-to-use-polishhook-symbol
\usepackage{amsmath}
\usepackage{amsfonts}
\usepackage{amssymb}
\usepackage{amsthm}
\usepackage{thmtools}
\usepackage{pgffor}
\usepackage{xcolor}
% \usepackage[lined,boxed,ruled,norelsize,algo2e,linesnumbered,noend]{algorithm2e}
\usepackage{cleveref}
\usepackage{graphicx}
\usepackage{geometry}
% \geometry{verbose,tmargin=1in,bmargin=1in,lmargin=1in,rmargin=1in}
\usepackage{enumitem}


% \newtheoremstyle{mystyle}
%   {12pt} % Space above
%   {24pt} % Space below
%   {} % Body font
%   {} % Indent amount
%   {\bfseries} % Theorem head font
%   {.} % Punctuation after theorem head
%   {.5em} % Space after theorem head
%   {} % Theorem head spec (can be left empty, meaning `normal')

% \theoremstyle{mystyle}


\newtheorem{assumption}{Assumption}
\newtheorem{definition}{Definition}
\newtheorem{theorem}{Theorem}
\newtheorem{lemma}{Lemma}
\newtheorem{proposition}{Proposition}
\newtheorem{corollary}{Corollary}

\usepackage{thmtools}

% \declaretheorem[refname={Theorem,Theorems},style=mystyle, Refname={Theorem,Theorems}]{theorem}
% \declaretheorem[numberlike=theorem, style=mystyle]{lemma}
% \declaretheorem[numberlike=theorem, style=mystyle]{invariant}
% \declaretheorem[numberlike=theorem, style=mystyle]{corollary}
% \declaretheorem[numberlike=theorem, style=mystyle]{definition}
% \declaretheorem[numberlike=theorem, style=mystyle]{claim}
% \declaretheorem[numberlike=theorem, style=mystyle]{fact}
% \declaretheorem[numberlike=theorem, style=mystyle]{assumption}
% \declaretheorem[numbered=no]{remark}

% number sets
\newcommand{\Real}{\mathbb{R}}
\newcommand{\Nat}{\mathbb{N}}
\newcommand{\F}{\mathbb{F}}
\newcommand{\Z}{\mathbb{Z}}
\renewcommand{\S}{\mathbb{S}}%

% probability
\renewcommand{\P}{\mathbb{P}}
\newcommand{\E}{\mathbb{E}}

% redfine some classical notation
\renewcommand{\tilde}{\widetilde}
\renewcommand{\hat}{\widehat}
\renewcommand{\bar}{\overline}

% operators
\newcommand{\median}{\operatorname{median}}

\DeclareMathOperator*{\argmin}{argmin}
\DeclareMathOperator*{\argmax}{argmax}
\newcommand{\poly}{\operatorname{poly}}
\newcommand{\polylog}{\operatorname{polylog}}
\newcommand{\mdiag}{\mathbf{Diag}} % diagonal matrix with vector on diag
\newcommand{\sign}{\operatorname{sign}}%
\newcommand{\nnz}{\operatorname{nnz}}%
\newcommand{\tail}{\operatorname{tail}}%
\newcommand{\dist}{\operatorname{dist}}%

% matrices
% \mX is \mathbf{X}
\foreach \x in {A,...,Z}{%
	\expandafter\xdef\csname m\x\endcsname{\noexpand\mathbf{\x}}
}
% \omX is \overline{\mathbf{X}}
\foreach \x in {A,...,Z}{%
	\expandafter\xdef\csname om\x\endcsname{\noexpand\overline{\noexpand\mathbf{\x}}}
}

\foreach \x in {A,...,Z}{%
	\expandafter\xdef\csname c\x\endcsname{\noexpand\mathcal{\x}}
}




\newcommand{\ob}{\overline{b}}
\newcommand{\ox}{\overline{x}}
\newcommand{\os}{\overline{s}}


\usepackage{mathtools}


% Reduce paragraph spacing https://tex.stackexchange.com/questions/4891/how-do-i-control-the-spacing-above-a-new-paragraph
\makeatletter
\renewcommand{\paragraph}{%
	\@startsection{paragraph}{4}%
	{\z@}{1.25ex \@plus 1ex \@minus .2ex}{-1em}%
	{\normalfont\normalsize\bfseries}%
}
\makeatother

\def\norm#1{\left\| #1 \right\|}



% \usepackage[
%     sorting=nyt,
% 	style=alphabetic,
%     maxalphanames=3,
%     minalphanames=3,
%     maxcitenames=3,
%     minnames=1,
%     maxbibnames=20,
% 	backref=true,
% 	doi=true,
% 	url=true,
% 	backend=biber
% ]{biblatex}
\usepackage{tcolorbox}


\newcommand{\thmbox}[1]{%
  \begin{tcolorbox}[
      colback = blue!5!white,
      colframe = blue!75!black,
      width          = \dimexpr\textwidth + 1cm\relax,
      % width    = \textwidth,   % full text width
      % left     = 5pt,          % no extra padding on the left
      % right    = 5pt,           % no extra padding on the right
      enlarge left by  = -0.5cm,
      enlarge right by = -0.5cm
  ]
    #1
  \end{tcolorbox}%
}

\newcommand{\mainthmbox}[1]{%
  \begin{tcolorbox}[colback=blue!5!white,colframe=blue!75!black,boxrule=3pt,
    ]
    #1
  \end{tcolorbox}%
}


\newcommand{\defbox}[1]{%
  \begin{tcolorbox}[colback=orange!5!white,colframe=orange!75!black, 
    width          = \dimexpr\textwidth + 1cm\relax,
    % width    = \textwidth,   % full text width
    % left     = 5pt,          % no extra padding on the left
    % right    = 5pt,           % no extra padding on the right
    enlarge left by  = -0.5cm,
    enlarge right by = -0.5cm]
    #1
  \end{tcolorbox}%
}

\newcommand{\blockbox}[1]{%
  \begin{tcolorbox}[colback=yellow!5!white,colframe=yellow!75!black]
    #1
  \end{tcolorbox}%
}

\newcommand{\shortminus}{\scalebox{0.5}[1.0]{\( - \)}}

% \usepackage{mathabx}  

\usepackage{cancel}
% \usepackage{arydshln}

\usepackage{geometry}
 \geometry{
        a4paper,
        total={170mm,257mm},
        left=20mm,
        right=10mm,
        top=20mm,
        }
        
% \usepackage{graphicx} % Required for inserting images

\usepackage{tikz}

\usetikzlibrary{arrows.meta,
                positioning,
                shadows}
                
\usetikzlibrary{shapes,snakes}
\usepackage{bm}
\usepackage{xspace}

\usepackage{cancel}


\setlist[itemize]{topsep=1pt, itemsep=1pt, parsep=1pt, leftmargin=*}

\usepackage{subcaption}
% \usepackage{graphicx}


  

 % \title{Transformers can Learn the Energy Spaces of Functions}
\title{Learning Mixtures of Anisotropic Gaussian Kernels with Transformers}
\newcommand{\energy}{{E}}
\newcommand{\nsamples}{{n}}
\newcommand{\nfunctions}{{m}}
\author{}

\begin{document}


\maketitle

% \newcommand{\manifold}{\mathcal{M}}
% \newcommand{\dataset}{{D}}

\section{Introduction}


\begin{abstract}
	Transformers with billions of parameters have achieved remarkable success across various domains, such as language modeling and computer vision.
	%
	Existing theoretical works analyze the empirically observed phenomena, such as training dynamics, in-context learning (ICL) generalization, depth/width effects, scaling laws, only in isolation and in restricted settings (single-layer, linear regression, fixed width), leaving open whether they are manifestations of a single underlying mechanism.
	%
	In this paper, we show how a Transformer can provably learn a Gaussian Process (GP) with covariance kernel $\kernel$ over a function space $p_\kernel(f) \propto \exp(-\frac{1}{2}\|f\|_{\HK}^2)$ from noisy observations $f_j(x_i) + \epsilon$.
	%
	We consider a kernel being mixture of anisotropic Gaussian kernels parametrised by a mixing measure $\theta$ and show how in the mean-field regime, where the number of heads $\to \infty$, training a Transformer recovers $\theta$ with provable approximation and generalization guarantees.
	%
	Our findings prove the convergence of the unsupervised training dynamics of Transformers, the emergence of scaling laws in the considered setting, and the necessity of the regularization of $\theta$ for the generalization properties of the ICL.
	%
	We evaluate our architecture on synthetic and real-world benchmarks to support our theoretical findings.
\end{abstract}

% \begin{theorem}[Main Result]
% 	\begin{align*}
% 		error = n^{-\alpha} + m^{-\beta} + H^{-\gamma} + D^{-\delta} + T^{-\tau}
% 	\end{align*}
% \end{theorem}

This paper makes the following contributions:
\begin{itemize}
	\item \textbf{Learning distributions over functions.}
	      We formulate the problem of learning a Gaussian distribution $\mathcal{GP}(0, \kernel)$ over a function space such that $p(f) \propto \exp(-\frac{1}{2}\|f\|_{\HK}^2)$, where $\kernel$ is an integral anisotropic Gaussian kernel,
	      and the function space $\ffam$ is the reproducing kernel Hilbert space $\HK$ associated with $\kernel$, from noisy observations as a regularized optimization problem.


	\item \textbf{Connection to in-context learning.}
	      We draw a connection of this problem to the in-context learning in Transformers Models, discuss how it can explain the convergence to the optimal solution, the empirically observed scaling laws, the role of the positional encodings and the importance of the regularization for the generalization properties.

	\item \textbf{Architectural implications and empirical validation.}
	      Guided by the optimization problem, we propose a minor modification of the Transformer
	      architecture that enforces the structural assumptions of our framework.
	      Experiments on both synthetic and real-world benchmarks demonstrate improvements
	      in size-generalization, interpretability, and computational efficiency.
\end{itemize}

\newpage

\section{Related Work}

Several works have showed that a fixed Transformer can implemet a gradient descent for a simple class of functions such as linear regression, Bayesian inference, etc.

\newpage


\subsection{Problem Setup}\label{sec:set-up}

We are given a domain $\domain$ equipped with a known probability measure $\mu$,
a function space $\ffam$ of functions $f \colon \domain \to \Real$,
endowed with probability measure $\nu(f) \propto \exp(-\frac{1}{2}\|f\|_{\HK}^2))$ where $\HK$ is the reproducing kernel Hilbert space of functions induced by a \emph{unknown} kernel $\kernel$.
We observe a dataset $\dataset = \left\{ D^{(1)}, \ldots, D^{(\nfunctions)} \right\}$ consisting of $\nfunctions$ observations $D^{(j)}$ of functions $f_1, \ldots, f_{\nfunctions} \sim \nu(\ffam)$
each evaluated at $\nsamples$ points drawn from $\mu$:
\begin{gather*}
	\mathbf{X} = \{X_1, \ldots, X_{\nsamples}\} \overset{\text{i.i.d.}}{\sim} \mu(\domain), \qquad
	\epsilon_{j,i} \overset{\text{i.i.d.}}{\sim} \mathcal{N}(0, \sigma^2), \\
	\dataset = \Bigl\{
	\bigl\{ (X_i,\, f_j(X_i) + \epsilon_{j,i}) \bigr\}_{i=1}^{\nsamples}
	\Bigr\}_{j=1}^{\nfunctions}.
\end{gather*}

Our goal is to learn the measure $\nu$ over $\ffam$.
Given $\nu$, predicting a new function $f \sim \nu$ from noisy observations
$\{(x_i, y_i)\}_{i=1}^{\nsamples}$, where $y_i = f(x_i) + \epsilon_i$ and
$\epsilon_i \overset{\text{i.i.d.}}{\sim} \mathcal{N}(0, \sigma^2)$,
reduces to MAP estimation:
\begin{align*}
	\hat{f}
	= \arg\min_{f \in \ffam}
	\sum_{i=1}^{\nsamples} \bigl| y_i - f(x_i) \bigr|^2
	+ \frac{1}{2}\|f\|_{\HK}^2
\end{align*}


\paragraph{Connection to ICL in Transformers.}

The setup described above is closely related to the in-context learning in Transformers.
Given a prompt $(x_1, y_1), \ldots, (x_{\nsamples}, y_{\nsamples})$, where $x_i$ are positional encodings
and $y_i$ are tokens, the task is to predict the next token at position $x_{\nsamples+1}$.
Each prompt can be viewed as a noisy observation of a latent function $f \sim \exp(-\frac{1}{2}\|f\|_{\HK}^2)$,
and distinct prompts correspond to distinct draws from $\exp(-\frac{1}{2}\|f\|_{\HK}^2)$.
The Transformer thus implicitly learns the shared structure of the functions,
that is, the kernel $\kernel$ itself, and uses it to generalize within a new context.

% \section{TODO}\label{sec:linear-energy-functionals}
\section{Gaussian Measures on the Reproducing Kernel Hilbert Space}\label{sec:gaussian-measure-on-the-rkhs}

In this section, we formally define the Gaussian Process in the Hilbert space $\nu(f) \propto \exp(-\frac{1}{2}\|f\|_{\HK}^2))$ and discuss its properties for anisotropic Gaussian kernels on the domain $[0, 1]^D$. We also show how knowing the kernel $\kernel$ and the observed noisy functions values $Y_i = f(X_i) + \varepsilon_i$ for $f \sim \mu_\kernel, i \in [n]$, provides a provable generalization guarantees for the kernel ridge regression estimator $\hat{f}_n$.

% \subsection{Assumptions}




%%%%%%%%%%%%%%%%%%%%%%%%%%%%%%%%%%%%%%%%%%%%%%%%%%%%%%%%%%%%%%%%%%%%%%%%%%%%%%%
% \section{Setup and Notation}
%%%%%%%%%%%%%%%%%%%%%%%%%%%%%%%%%%%%%%%%%%%%%%%%%%%%%%%%%%%%%%%%%%%%%%%%%%%%%%%

Let $(\X, \mathcal{A}, P_X)$ be a probability space where $\X$ is the input domain. Let $\kernel: \X \times \X \to \R$ be a symmetric positive semi-definite (PSD) kernel satisfying the integrability condition:
$
	\int_{\X} \kernel(x, x) \, dP_X(x) < \infty.
$
We work primarily in the Hilbert space $L^2(P_X)$ with inner product
$
	\inprod{f}{g}_{L^2} = \int_{\X} f(x) g(x) \, dP_X(x).
$

%%%%%%%%%%%%%%%%%%%%%%%%%%%%%%%%%%%%%%%%%%%%%%%%%%%%%%%%%%%%%%%%%%%%%%%%%%%%%%%
% \section{The Integral Operator}
%%%%%%%%%%%%%%%%%%%%%%%%%%%%%%%%%%%%%%%%%%%%%%%%%%%%%%%%%%%%%%%%%%%%%%%%%%%%%%%

\begin{definition}[Integral Operator]
	The integral operator $T_\kernel: L^2(P_X) \to L^2(P_X)$ associated with kernel $\kernel$ is defined by
	\begin{equation}
		(T_\kernel g)(x) = \int_{\X} \kernel(x, y) \, g(y) \, dP_X(y).
	\end{equation}
\end{definition}

\begin{theorem}[Spectral Decomposition]
	\label{thm:spectral}
	The operator $T_\kernel$ is self-adjoint, positive, and trace-class. It admits a spectral decomposition
	$
		T_\kernel \phi_k = \mu_\kernel \phi_k, k \in \N,
	$
	where:
	\begin{enumerate}[label=(\roman*)]
		\item $\mu_1 \geq \mu_2 \geq \cdots > 0$ are the positive eigenvalues (counted with multiplicity),
		\item $\{\phi_k\}_{k=1}^\infty$ is an orthonormal basis of $\overline{\Range(T_\kernel)} \subseteq L^2(P_X)$,
		\item $\sum_{k=1}^\infty \mu_\kernel = \Tr(T_\kernel) = \int_{\X} K(x,x) \, dP_X(x) < \infty$.
	\end{enumerate}
	The kernel admits the Mercer representation
	$
		K(x, y) = \sum_{k=1}^\infty \mu_\kernel \, \phi_k(x) \, \phi_k(y),
	$
	with convergence in $L^2(P_X \otimes P_X)$.
\end{theorem}

%%%%%%%%%%%%%%%%%%%%%%%%%%%%%%%%%%%%%%%%%%%%%%%%%%%%%%%%%%%%%%%%%%%%%%%%%%%%%%%
% \section{The Reproducing Kernel Hilbert Space}
%%%%%%%%%%%%%%%%%%%%%%%%%%%%%%%%%%%%%%%%%%%%%%%%%%%%%%%%%%%%%%%%%%%%%%%%%%%%%%%

\begin{definition}[RKHS via Spectral Representation]
	\label{def:rkhs}
	The reproducing kernel Hilbert space (RKHS) $\HK$ associated with $\kernel$ is
	\begin{equation}
		\HK = \left\{ f = \sum_{k=1}^\infty \theta_k \phi_k : \norm{f}_{\HK}^2 := \sum_{k=1}^\infty \frac{\theta_k^2}{\mu_\kernel} < \infty \right\},
	\end{equation}
	equipped with the inner product
	\begin{equation}
		\inprod{f}{g}_{\HK} = \sum_{k=1}^\infty \frac{\theta_k \eta_k}{\mu_\kernel}, \quad \text{where } f = \sum_k \theta_k \phi_k, \; g = \sum_k \eta_k \phi_k.
	\end{equation}
\end{definition}

\begin{proposition}[Operator Characterization]
	\label{prop:rkhs-operator}
	The RKHS admits the equivalent characterization
	$
		\HK = \Range(T_\kernel^{1/2})
	$
	with norm
	$
		\norm{f}_{\HK} = \norm{T_\kernel^{-1/2} f}_{L^2}.
	$
	The inclusion $\HK \hookrightarrow L^2(P_X)$ is continuous and dense in $\overline{\Range(T_\kernel)}$.
\end{proposition}

\begin{remark}
	The space $\HK$ is ``smaller'' than $L^2(P_X)$: for $f \in \HK$, the coefficients $\theta_k = \inprod{f}{\phi_k}_{L^2}$ must decay faster than $\sqrt{\mu_\kernel}$.
\end{remark}

%%%%%%%%%%%%%%%%%%%%%%%%%%%%%%%%%%%%%%%%%%%%%%%%%%%%%%%%%%%%%%%%%%%%%%%%%%%%%%%
% \section{The Gaussian Measure}
%%%%%%%%%%%%%%%%%%%%%%%%%%%%%%%%%%%%%%%%%%%%%%%%%%%%%%%%%%%%%%%%%%%%%%%%%%%%%%%

\begin{definition}[Gaussian Measure via Characteristic Functional]
	\label{def:gaussian-measure}
	The \emph{Gaussian measure} $\mu_\kernel$ on $L^2(P_X)$ is the unique Borel probability measure with characteristic functional
	\begin{equation}
		% \boxed{
		\hat{\mu}_K(g) := \int_{L^2(P_X)} e^{i\inprod{g}{f}_{L^2}} \, d\mu_\kernel(f) = \exp\left( -\frac{1}{2} \inprod{g}{T_\kernel g}_{L^2} \right)
		% }
	\end{equation}
	for all $g \in L^2(P_X)$.
\end{definition}

\begin{theorem}[Existence and Characterization]
	\label{thm:gaussian-existence}
	The measure $\mu_\kernel$ exists, is unique, and satisfies the following:

	\begin{enumerate}[label=(\roman*)]
		\item \textbf{Karhunen-Lo\`eve representation:}
		      \begin{equation}
			      %   \boxed{
			      f \sim \mu_\kernel \iff f = \sum_{k=1}^\infty \sqrt{\mu_\kernel} \, \xi_k \, \phi_k, \quad \xi_k \iid N(0,1)
			      %   }
		      \end{equation}
		      where the series converges in $L^2(\Omega; L^2(P_X))$.

		\item \textbf{Mean and covariance:}
		      \begin{align}
			      \E_{\mu_\kernel}[f]                                       & = 0,                                                              \\
			      \E_{\mu_\kernel}[\inprod{f}{g}_{L^2} \inprod{f}{h}_{L^2}] & = \inprod{g}{T_\kernel h}_{L^2}, \quad \forall g, h \in L^2(P_X).
		      \end{align}
		      Equivalently, $\E_{\mu_\kernel}[f \otimes f] = T_\kernel$ as operators.

		\item \textbf{Finite-dimensional distributions:} For any $x_1, \ldots, x_n \in \X$,
		      \begin{equation}
			      (f(x_1), \ldots, f(x_n))^\top \sim N(0, K_{\mathbf{X}\mathbf{X}}),
		      \end{equation}
		      where $(K_{\mathbf{X}\mathbf{X}})_{ij} = K(x_i, x_j)$.
	\end{enumerate}
\end{theorem}

\begin{remark}[Gaussian Process Equivalence]
	The measure $\mu_\kernel$ is the law of the Gaussian process $f \sim \mathcal{GP}(0, K)$.
\end{remark}

%%%%%%%%%%%%%%%%%%%%%%%%%%%%%%%%%%%%%%%%%%%%%%%%%%%%%%%%%%%%%%%%%%%%%%%%%%%%%%%
% \section{The Cameron-Martin Theorem}
%%%%%%%%%%%%%%%%%%%%%%%%%%%%%%%%%%%%%%%%%%%%%%%%%%%%%%%%%%%%%%%%%%%%%%%%%%%%%%%

The following theorem establishes $\HK$ as the Cameron-Martin space of $\mu_\kernel$, providing the fundamental link between the RKHS and the Gaussian measure.

\begin{theorem}[Cameron-Martin]
	\label{thm:cameron-martin}
	Let $\mu_\kernel$ be the Gaussian measure on $L^2(P_X)$ and $\HK$ its associated RKHS. Then:

	\begin{enumerate}[label=(\roman*)]
		\item \textbf{Null RKHS:} Sample paths lie outside the RKHS almost surely:
		      $
			      \mu_\kernel(\HK) = 0.
		      $

		\item \textbf{Admissible shifts:} For $h \in L^2(P_X)$, define the translated measure $\mu_\kernel^h(A) := \mu_\kernel(A - h)$. Then
		      \begin{equation}
			      \mu_\kernel^h \ll \mu_\kernel \iff h \in \HK.
		      \end{equation}
		      That is, $\HK$ is exactly the set of admissible shifts preserving absolute continuity.

		\item \textbf{Radon-Nikodym derivative:} For $h \in \HK$,
		      \begin{equation}
			      %   \boxed{
			      \frac{d\mu_\kernel^h}{d\mu_\kernel}(f) = \exp\left( \inprod{h}{f}_{\HK} - \frac{1}{2}\norm{h}_{\HK}^2 \right)
			      %   }
		      \end{equation}
		      where the pairing $\inprod{h}{f}_{\HK}$ extends to $f \in L^2(P_X)$ via
		      \begin{equation}
			      \inprod{h}{f}_{\HK} = \sum_{k=1}^\infty \frac{\eta_k \theta_k}{\mu_\kernel}, \quad h = \sum_k \eta_k \phi_k, \; f = \sum_k \theta_k \phi_k.
		      \end{equation}
	\end{enumerate}
\end{theorem}

\begin{corollary}
	The RKHS norm $\norm{\cdot}_{\HK}$ measures the ``cost'' of deterministic shifts under the Gaussian measure.
\end{corollary}

%%%%%%%%%%%%%%%%%%%%%%%%%%%%%%%%%%%%%%%%%%%%%%%%%%%%%%%%%%%%%%%%%%%%%%%%%%%%%%%
% \section{The Onsager-Machlup Functional}
%%%%%%%%%%%%%%%%%%%%%%%%%%%%%%%%%%%%%%%%%%%%%%%%%%%%%%%%%%%%%%%%%%%%%%%%%%%%%%%

The formal expression ``$p(f) \propto \exp(-\frac{1}{2}\norm{f}_{\HK}^2)$'' is made rigorous by the following theorem.

\begin{theorem}[Onsager-Machlup]
	\label{thm:onsager-machlup}
	Let $\mu_\kernel$ be a Gaussian measure on a Banach space $\B \supseteq \HK$ with Cameron-Martin space $\HK$. Let $\norm{\cdot}_{\B}$ be a measurable norm satisfying appropriate regularity conditions.\footnote{Specifically, $\norm{\cdot}_{\B}$ should be measurable and the small ball probabilities $\mu_\kernel(B_\epsilon(0))$ should satisfy certain asymptotic conditions; see Zeitouni (1989).} Then for any $f, g \in \HK$:
	\begin{equation}
		% \boxed{
		\lim_{\epsilon \to 0} \frac{\mu_\kernel\bigl(\{h \in \B : \norm{h - f}_{\B} < \epsilon\}\bigr)}{\mu_\kernel\bigl(\{h \in \B : \norm{h - g}_{\B} < \epsilon\}\bigr)} = \exp\left( -\frac{1}{2}\norm{f}_{\HK}^2 + \frac{1}{2}\norm{g}_{\HK}^2 \right)
		% }
	\end{equation}
\end{theorem}

\begin{corollary}[RKHS Norm as Log-Density]
	Taking $g = 0$:
	\begin{equation}
		\lim_{\epsilon \to 0} \frac{\mu_\kernel(B_\epsilon(f))}{\mu_\kernel(B_\epsilon(0))} = \exp\left( -\frac{1}{2}\norm{f}_{\HK}^2 \right).
	\end{equation}
	This is the precise sense in which
	$
		``p(f) \propto \exp\left(-\frac{1}{2}\norm{f}_{\HK}^2\right)."
	$
\end{corollary}

\begin{remark}
	The function $f \mapsto \frac{1}{2}\norm{f}_{\HK}^2$ is called the \emph{Onsager-Machlup functional} of $\mu_\kernel$. It plays the role of a ``negative log-density'' for the Gaussian measure.
\end{remark}

%%%%%%%%%%%%%%%%%%%%%%%%%%%%%%%%%%%%%%%%%%%%%%%%%%%%%%%%%%%%%%%%%%%%%%%%%%%%%%%
\subsection{Connection to Bayesian Inference and Kernel Ridge Regression}
%%%%%%%%%%%%%%%%%%%%%%%%%%%%%%%%%%%%%%%%%%%%%%%%%%%%%%%%%%%%%%%%%%%%%%%%%%%%%%%

The following theorem establishes the connection between the Bayesian inference and the kernel ridge regression estimator: knowing the kernel $\kernel$ and the observed noisy functions values $\mathbf{Y} = \{Y_i = f(X_i) + \varepsilon_i\}_{i=1}^n$ for $f \sim \mu_\kernel, \varepsilon_i \iid N(0, \sigma^2)$, how well can we estimate this unknown function $f$?

\begin{theorem}[Posterior and MAP Estimator]
	\label{thm:posterior}
	Suppose we observe
	\begin{equation}
		Y_i = f(X_i) + \varepsilon_i, \quad i = 1, \ldots, n,
	\end{equation}
	where $\varepsilon_i \iid N(0, \sigma^2)$, $X_i \iid P_X$, and the prior is $f \sim \mu_\kernel$. Then:

	\begin{enumerate}[label=(\roman*)]
		\item The posterior $\mu_\kernel^{(n)}(\cdot \mid \mathbf{Y}, \mathbf{X})$ is a Gaussian measure.

		\item The posterior mean coincides with the kernel ridge regression estimator:
		      \begin{equation}
			      %   \boxed{
			      \hat{f}_n = \argmin_{f \in \HK} \left\{ \frac{1}{\sigma^2} \sum_{i=1}^n (Y_i - f(X_i))^2 + \norm{f}_{\HK}^2 \right\}
			      %   }
		      \end{equation}

		\item In closed form:
		      \begin{equation}
			      \hat{f}_n(\cdot) = \mathbf{k}(\cdot)^\top (\kernel_{\mathbf{X}\mathbf{X}} + \sigma^2 I)^{-1} \mathbf{Y},
		      \end{equation}
		      where $\mathbf{k}(x) = (\kernel(x, X_1), \ldots, \kernel(x, X_n))^\top$.

		\item The generalization error of the kernel ridge regression estimator $\hat{f}_n$ is given by:
		      \begin{equation}
			      \norm{f - \hat{f}_n}_{\HK}^2  \leq TODO
		      \end{equation}
	\end{enumerate}

\end{theorem}

\begin{remark}
	The RKHS norm appears as the regularizer in the MAP estimator precisely because it is the Onsager-Machlup functional of the GP prior. This establishes the equivalence:
	\begin{center}
		\emph{GP regression with covariance $K$} $\longleftrightarrow$ \emph{Kernel ridge regression with RKHS $\HK$}.
	\end{center}
\end{remark}



\section{The Integral Anisotropic Gaussian Kernel}\label{sec:integral-anisotropic-gaussian-kernel}

Here we introduce a family of kernels considered in this work: integral anisotropic Gaussian kernels. They cover a wide range of kernels, including the isotropic Gaussian kernel, the radial Gaussian kernel, the finite mixture of anisotropic Gaussian kernels, the Wishart mixing measure, the rank-one (ridge function) measure, and the integral of the anisotropic Gaussian kernel.

\begin{definition}[Integral Anisotropic Gaussian Kernel]
	\label{def:integral-kernel}
	Let $\theta \in \mathcal{M}_+(\Splus)$ be a finite non-negative Borel measure on the cone of positive semi-definite matrices. The \emph{integral anisotropic Gaussian kernel} induced by $\theta$ is the function $K_\theta: \R^d \times \R^d \to \R$ defined by
	\begin{equation}
		% \boxed{
		K_\theta(x, y) = \int_{A \in \Splus} K_A(x, y) \, d\theta(A) = \int_{A \in \Splus} \exp\left( -\frac{1}{2}(x-y)^\top A (x-y) \right) d\theta(A).
		% }
	\end{equation}
	We call $\theta$ the \emph{mixing measure} or \emph{structure measure} of the kernel.
\end{definition}



Given functions $f_1, \ldots, f_{\nfunctions} \sim \mu_{\kernel_\theta}$, where $\kernel_\theta$ is an integral anisotropic Gaussian kernel the MLE estimator of correponding mixing measure $\theta$ is given by:
\begin{equation}\label{eq:mle-estimator}
	\hat{\theta}_m = \arg\max_{\theta \in \Theta} \left\{ \prod_{i=1}^{\nfunctions} p(f_i \mid \kernel_\theta) \right\}
\end{equation}
where $p(f_i \mid \kernel_\theta) = \exp(-\frac{1}{2}\|f_i\|_{\HK}^2)$. Depending on the $\theta^*$ we can get the error rate of its MLE estimator $\hat{\theta}_m$, and the simpler the $\theta^*$ is, the better the generalization guarantees are. The following notion defines the simplicity of the mixing measure $\theta^*$ as its effective dimension.


\begin{definition}[Effective Dimension]
	\label{def:eff-dim}
	The \textbf{effective dimension} of $\theta^*$ relative to the model $\Theta$ is:
	\begin{equation}
		D(\theta^*) := \lim_{\epsilon \to 0} \frac{\log N(\epsilon, \Theta_{\theta^*}, d_K)}{\log(1/\epsilon)}
	\end{equation}
	where $\Theta_{\theta^*}$ is a local neighborhood of $\theta^*$ in $\Theta$, and $N(\epsilon, \cdot, d_K)$ is the $\epsilon$-covering number in the kernel metric $d_K(\theta_1, \theta_2) = \|K_{\theta_1} - K_{\theta_2}\|$.
\end{definition}

\begin{remark}
	The effective dimension $D(\theta^*)$ captures the \emph{local complexity} of the model around $\theta^*$. It can vary with $\theta^*$: the simpler the $\theta^*$ is, the smaller the effective dimension is, the better the generalization guarantees are.
\end{remark}



% \begin{tcolorbox}[colback=blue!5!white, colframe=blue!50!black, title=Main Result]
\begin{theorem}[General Rate Theorem]
	\label{thm:theta-mle-rate}
	Let $\theta^* \in \Theta \subseteq \mathcal{M}_+(\Splus)$. Suppose the model $\{\kernel_\theta : \theta \in \Theta\}$ has \emph{local metric entropy} around $\theta^*$ satisfying:
	\begin{equation}
		\log N\left(\epsilon, \{\theta \in \Theta : d_K(\kernel_\theta, \kernel_{\theta^*}) \leq \delta\}, d_K\right) \leq D(\theta^*) \cdot \log\left(\frac{\delta}{\epsilon}\right)
	\end{equation}
	for all $0 < \epsilon < \delta$ sufficiently small. Then:
	\begin{equation}
		% \boxed{
		d_K(\hat{\theta}_m, \theta^*) = O_p\left(m^{-\frac{1}{2 + D(\theta^*)/s}}\right)
		% }
	\end{equation}
	where $s$ is the smoothness parameter if $\theta^*$ has a density, or $s = \infty$ (giving rate $m^{-1/2}$) in the parametric case.

	More precisely:
	\begin{equation}
		% \boxed{
		d_K(\hat{\theta}_m, \theta^*) = O_p\left(m^{-\alpha(\theta^*)}\right), \quad \alpha(\theta^*) = \frac{1}{2 + D(\theta^*)}
		% }
	\end{equation}
	when the density has regularity $s = 1$.
\end{theorem}
% \end{tcolorbox}





% \begin{table}[h]
% 	\centering
% 	\caption{Notation}
% 	\label{tab:notation}
% 	\begin{tabular}{ll}
% 		\toprule
% 		\textbf{Symbol}                      & \textbf{Description}                                                       \\
% 		\midrule
% 		$\domain$                            & Domain, $\domain \subset \mathbb{R}^d$                                     \\
% 		$d$                                  & Ambient dimension                                                          \\
% 		$\mathbf{S}_+^d$                     & Set of $d \times d$ symmetric positive semi-definite matrices              \\
% 		$\mathcal{Q} \subset \mathbf{S}_+^d$ & Admissible set of metric tensors                                           \\
% 		$Q \in \mathcal{Q}$                  & PSD matrix defining an anisotropic metric                                  \\
% 		$\Delta_Q$                           & Anisotropic Laplacian operator, $\Delta_Q f := \nabla \cdot (Q \nabla f)$  \\
% 		$\mathcal{M}(\domain)$               & Space of Radon measures on $\domain$                                       \\
% 		$\|\cdot\|_{\mathcal{M}}$            & Total variation norm, $\|g\|_{\mathcal{M}} := \int_{\domain} |g(x)| \, dx$ \\
% 		\bottomrule
% 	\end{tabular}
% \end{table}


% \subsection{Measure-Based and Neural Representations}

% \paragraph{Example 8 (Barron / ReLU representation).}
% Let
% \[
% 	f_\mu(x) = \int_{\mathbb S^{d-1}} \sigma(\langle w,x\rangle)\, d\mu(w),
% \]
% with $\mu$ a finite signed Radon measure.
% Define
% \[
% 	\Phi : \mathcal B \to \mathcal M_+(\mathbb S^{d-1}),
% 	\qquad
% 	\Phi(f_\mu) := |\mu|.
% \]

% Here:
% \begin{itemize}
% 	\item domain: Barron space $\mathcal B$,
% 	\item codomain: finite nonnegative measures,
% 	\item regularizer: $R(f_\mu)=\|\mu\|_{\mathrm{TV}}$.
% \end{itemize}

% This promotes sparse ridge representations.


% \paragraph{Example 9 (Finsler--TV / Radon--Laplacian).}
% Combining the previous examples,
% \[
% 	\Phi : \mathcal B \to \mathcal C(\mathcal S_+,\mathbb{R}_+),
% 	\qquad
% 	\Phi(f_\mu)(Q) := \int \sqrt{w^\top Q w}\,|\mu|(dw).
% \]

% Equivalently,
% \[
% 	R(f_\mu)
% 	=
% 	\int_{\mathcal S_+} \Phi(f_\mu)(Q)\, d\theta(Q)
% 	=
% 	\int F(w)\,|\mu|(dw),
% 	\quad
% 	F(w):=\int \sqrt{w^\top Q w}\, d\theta(Q).
% \]

% This is a Finsler total variation norm on kink directions.


% \subsection{How transformer forward pass can implement a (sub)gradient descent}

% Now lets go through some of the examples in \Cref{par:anisotropic-second-order-tv}
% to see how transformer-like networks can implement a (sub)gradient $g$ descent for the inner problem.

% \paragraph{Energy regularization}

% \[
% 	\Phi : L^2(M) \to \mathbb{R},
% 	\qquad
% 	\Phi(f) := \|f\|_{L^2(M)}^2 \qquad R_\theta(f) = \theta \|f\|_{L^2(M)}^2
% 	\qquad
% 	g = 2f
% \]



% \paragraph{Anisotropic second-order TV}

% \[
% 	\Phi : X \to \mathcal C(\mathcal S_+,\mathbb{R}_+),
% 	\qquad
% 	\Phi(f)(Q) := \|\Delta_Q f\|_{\mathcal M},
% 	\qquad
% 	g = \int_{\mathcal S_+}  \Delta_Q \sign(\Delta_Q f) d\theta(Q)
% \]
% where $\Delta_Q f := \nabla\cdot(Q\nabla f)$.


% As the functions are defined on a "flat" domain $\domain \subset R^{d}$, the discrete version of $\Delta_Q$ looks like:
% \begin{align*}
% 	L \in R^{n \times n}: L_{ij} = \exp\left(\frac{-\|x_i - x_j\|_{Q^{-1}}^2}{\epsilon} \right)
% \end{align*}


% \subsection{Function Extrapolation}

% First, let us see how knowing the energy functional enables function extrapolation.
% Suppose we observe $D = \{(x_i, y_i)\}_{i=1}^{\nsamples}$ with $y_i = f(x_i) + \epsilon_i$,
% $\epsilon_i \overset{\text{i.i.d.}}{\sim} \mathcal{N}(0, \sigma^2)$, and wish to predict
% $y_{m+1} = f(x_{m+1})$ at a new point $x_{m+1}$.
% Knowing the Gibbs prior $p(f) \propto \exp(-\energy_{\theta^*}(f)) = \exp(-\langle \theta^*, \Phi(f) \rangle)$,
% we recover $f$ via MAP estimation:
% \begin{align*}
% 	\hat{f}_{\nsamples}
% 	 & = \argmax_{f \in \ffam}\;
% 	% p(f \mid \dataset_m)                                \\
% 	%  & = \argmax_{f \in \ffam}\;
% 	p\left(D \mid f\right)\, p(f)
% 	= \argmin_{f \in \ffam}\;
% 	\frac{1}{\sigma^2}\sum_{i=1}^{\nsamples} \bigl|y_i - f(x_i)\bigr|^2
% 	+ \energy_{\theta^*}(f).
% \end{align*}
% This estimation will have a generalization guarantee of extrapolating $f$ to a new point $x_{\nsamples+1} \sim \mu$ via the following theorem:
% \begin{theorem}[]\label{thm:f-generalization}
% 	\begin{align*}
% 		\left\|f^*-\hat{f}_{\nsamples}\right\|_{L^2} \lesssim \underbrace{\left\|f^*-f_{\theta^*}\right\|_{L^2}}_{\text{bias}}+\sigma \sqrt{\frac{\deff\left(\theta^*\right)}{\nsamples}}
% 	\end{align*}

% 	Correspondingly taking an expectation over $f \sim p_{\theta^*}$ yields:
% 	\begin{align*}
% 		\E_{f \sim p_{\theta^*}}\left[ \left\|f-\hat{f}_{\nsamples}\right\|_{L^2} \right] \lesssim \underbrace{\left\|f-f_{\theta^*}\right\|_{L^2}}_{\text{bias}}+\sigma \sqrt{\frac{\deff\left(\theta^*\right)}{\nsamples}}
% 	\end{align*}
% \end{theorem}

% For example, assuming that our function is $s$-smooth on $\domain \subset \Real^D, \dim(\domain) = d$ then \Cref{thm:f-generalization} recovers the standard rate ${O}\left( \nsamples^{-2s/(2s+d)} \right)$ of the minimax rate of estimation of a smooth function from $\nsamples$ noisy observations.


\section{Making the Optimization Problem Tractable via Transformers}\label{sec:optimization-problem}


Even though, the \Cref{thm:theta-mle-rate} provides a provable estimation error rate guarantees for recovering the mixing measure $\theta^*$, hence, the kernel $\kernel_\theta$, from the functions $f_1, \ldots, f_{\nfunctions} \sim \mu_{\kernel_\theta}$, in practice we only observe the noisy functions values $Y = f_j(X) + \varepsilon_i$ for $f_j \sim \mu_\kernel, \varepsilon_i \iid N(0, \sigma^2)$.
% Hence, we need to learn the mixing measure $\theta^*$ from the noisy functions values $Y_i = f_j(X_i) + \varepsilon_i$.
So given observations $\dataset = \left\{\left(\mathbf{X}^{(i)}, \mathbf{y}^{(i)}\right)\right\}_{i=1}^{\nfunctions}$ our optimization problem turns into
\begin{gather}\label{eq:loss-w-data}
	\hat \theta_m = \argmax_{\theta \in \Theta} p_\theta(\dataset) =
	\argmin_{\theta \in \Theta}
	-\frac{1}{\nfunctions} \sum_{i=1}^{\nfunctions}
	\log p_\theta\left(\mathbf{y}^{(i)} \mid \mathbf{X}^{(i)}\right)
\end{gather}
where
\begin{align*}
	- \log p_\theta(\mathbf{y} \mid \mathbf{X}) =
	-\log \int \exp\left(-\frac{1}{2\sigma^2}\|\mathbf{y} - f(\mathbf{X})\|^2\right) d\mu_{\kernel_\theta}(f)
\end{align*}

As our prior is a Gaussian Process, the exact Laplace approximation yields a closed-form expression for the log-likelihood:
\begin{align*}
	- \log p_\theta(\mathbf{y} \mid \mathbf{X}) =
	% L_\theta(\hat{f}_\theta, D)
	% +
	% \frac{1}{2}\log \det H_\theta
	% +
	% \frac{n}{2}\log(2\pi)
	% \\
	% =
	\frac{1}{2} \mathbf{y}^\top (\kernel_\theta(\mathbf{X}, \mathbf{X}) + \sigma^2 I)^{-1} \mathbf{y}
	+ \frac{1}{2} \log \det(\kernel_\theta(\mathbf{X}, \mathbf{X}) + \sigma^2 I)
	+ \frac{\nsamples}{2} \log(2\pi)
\end{align*}

Even though this problem is convex, there are two main challenges in implementing it in practice via gradient descent:
\begin{itemize}
	\item Parameter $\theta$ is infinite-dimensional.
	\item The matrix inversion $(\kernel_\theta(\mathbf{X}, \mathbf{X}) + \sigma^2 I)^{-1}$ and the gradient with respect to $\theta$ is computationally prohibitive for large $\nsamples$.
\end{itemize}

First, let us discuss how to avoid the matrix inversion.

\subsection{Bilevel Optimization Via Masking}\label{sec:masking}
One way to avoid the matrix inversion is to turn the optimization problem \eqref{eq:loss-w-data} into a bilevel optimization problem via masking approach.
\begin{align*}
	\frac{1}{2} \mathbf{y}^\top (\kernel_\theta + \sigma^2 I)^{-1} \mathbf{y}
	= \frac{1}{2\sigma^2} \|\mathbf{y} - \hat{f}_\theta(\mathbf{X})\|^2
	+
	\frac{1}{2} \| f \|_{\kernel_\theta}^2
\end{align*}
where via \Cref{thm:posterior} we have
\begin{align*}
	\hat{f}_\theta =
	\argmin_{f \in \ffam}
	\frac{1}{2\sigma^2} \|\mathbf{y} - f(\mathbf{X})\|_2^2 + \frac{1}{2} \| f \|_{\kernel_\theta}^2
\end{align*}


Given the collection of datasets $\mathcal{D} = \left\{\left(\mathbf{X}^{(i)}, \mathbf{y}^{(i)}\right)\right\}_{i=1}^{\nfunctions}$ we can split each pair $\left(\mathbf{X}, \mathbf{y}\right)$ into a nonoverlapping masked and unmasked halves: $\left(\mathbf{X}, \mathbf{y}\right) = \left(\mathbf{X}_{\text{C}}, \mathbf{y}_{\text{C}}\right) \bigsqcup \left(\mathbf{X}_{\text{Q}}, \mathbf{y}_{\text{Q}}\right)$.

Substituting this into \Cref{eq:loss-w-data} we have
\begin{gather*}
	\hat \theta_{\dataset} =
	\min_{\theta \in \Theta}
	\frac{1}{\nfunctions} \sum_{i=1}^\nfunctions
	\frac{1}{2\sigma^2} \|\mathbf{y}_{\text{Q}} - \hat{f}_\theta(\mathbf{X}_{\text{Q}})\|^2
	+
	\frac{1}{2} \| \hat{f}_\theta \|_{\kernel_\theta}^2
	+
	\frac{1}{2} \log \det(\kernel_\theta(\mathbf{X}_{\text{Q}}, \mathbf{X}_{\text{Q}}) + \sigma^2 I)
	% \quad \text{(outer problem)}
	\nonumber
	\\
	\text{where }
	\hat{f}_\theta =
	\argmin_{f \in \HK}
	\frac{1}{2\sigma^2} \|\mathbf{y}_{\text{C}} - f(\mathbf{X}_{\text{C}})\|_2^2 + \frac{1}{2} \| f \|_{\kernel_\theta}^2
\end{gather*}

% This is bilever optimization. Even though it lost the convexity, it is tractable and can be solved in practice by alternating gradient descent. 
The inner problem has parameter $\theta$ fixed and find the MLE of the function $\hat f_\theta$ given the unmasked dataset $\mathbf{X}_{\text{C}}$ and $\mathbf{y}_{\text{C}}$. The outer problem then finds the parameter $\theta$ that yields the best prediction of $\hat f_\theta$ on the masked dataset $\mathbf{X}_{\text{Q}}$ and $\mathbf{y}_{\text{Q}}$. Both problems are convex on their own but not jointly convex. However, we can solve it by alternating gradient descent with a following guarantee:
\begin{theorem}
	\todo{add theorem}
\end{theorem}

\paragraph{The Inner Problem.}
For a dataset $\mathbf{X}_{\text{C}}, \mathbf{y}_{\text{C}}$ we have
\begin{align*}
	\hat{f}_\theta =
	\argmin_{f \in \HK}
	\frac{1}{2\sigma^2} \|\mathbf{y}_{\text{C}} - f(\mathbf{X}_{\text{C}})\|_2^2 + \frac{1}{2} \| f \|_{\kernel_\theta}^2
\end{align*}

One gradient step of the inner problem is:

\begin{align*}
	f^{(t+1)} & = f^{(t)} - \eta \nabla_f \left[\frac{1}{2\sigma^2} \|\mathbf{y}_C - f(\mathbf{X}_C)\|_2^2 + \frac{1}{2}\|f\|_{\mathcal{H}_{\kernel_\theta}}^2\right]\bigg|_{f = f^{(t)}} \\
	          & = f^{(t)} - \eta \left( -\frac{1}{\sigma^2}\sum_{i=1}^{n_C} (y_i - f^{(t)}(x_i)) \kernel_\theta(\cdot, x_i) + f^{(t)} \right)                                             \\
	          & = (1 - \eta) f^{(t)} + \frac{\eta}{\sigma^2} \sum_{i=1}^{n_C} (y_i - f^{(t)}(x_i)) \kernel_\theta(\cdot, x_i)
\end{align*}


% Note, in all the examples that follow $\Phi(f)$ is convex in $f$.

\paragraph{The Outer Problem.}




% \section{Connection to Transformers And Its Implications}

% In most of the examples above, the key step of implementing a (sub)gradient descent for the inner problem is to compute the discrete (anisotropic) Laplacian $\Delta_Q$ for some PSD matrix $Q$. The original Laplacian $\Delta_Q$ can be approximated by its discrete version $\hat \Delta_Q \in \Real^{n \times n}$ computed on a set of points $x_1, \ldots, x_n \in \domain$:
% \begin{gather*}
% 	% \label{eq:discrete-laplacian}
% 	\hat \Delta_Q = I - D^{-\nicefrac{1}{2}} W D^{-\nicefrac{1}{2}} \quad \text{where} \\
% 	W_{ij} = \exp\left(\frac{-\langle x_i - x_j, \Pi_{x_i} Q^\dagger \Pi_{x_i} (x_i - x_j)\rangle}{\epsilon} \right) \quad \text{and} \quad D_{ii} = \sum_{j=1}^n W_{ij}.
% \end{gather*}

% Even though the attention mechanism is not exactly the same as the discrete Laplacian, it is similar to it in a few ways. The attention mechanism is a way to compute the weighted average ..... The discrete Laplacian is a way to compute the weighted average of the features of the points in the neighborhood of a given point.
% Consider a transformer-like architecture where the softmax attention is replaced with 2 different types of heads per layer:
% \begin{itemize}
% 	\item 1 \textbf{linear} head: $g_0(y^{(t)}) = \ldots$
% 	\item $H$ \textbf{diffusion} heads each parametrized by some $Q_{h} \in \mathcal{Q}$ and $\theta_{h} \in \Real^+$: $\hat g_i(y^{(t)}) =  \eta \theta_h \Delta_{Q_h} y^{(t)}$.
% \end{itemize}
% Then the forward pass is
% \begin{align*}
% 	y^{(t+1)} \in \Real^{\nsamples}, \hat g_i:  \Real^{\nsamples} \to \Real^{\nsamples}:
% 	% \\
% 	% \quad \quad
% 	y^{(t+1)} = \left(I - \left[y^{(t)} - y^{(0)}\right]_\text{mask} - \sum_{h=1}^H \hat g_h(y^{(t)})\right) y^{(t)}
% \end{align*}
% Which in the {mean-field limit} $\nsamples \to \infty, H \to \infty$ yields
% \begin{align*}
% 	f^{(t+1)} \in \ffam: g_i: \ffam \to \ffam:
% 	f^{(t+1)} = \left(I - \left[f^{(t)} - f^{(0)}\right]_\text{mask} - \int\limits_{Q\in\mathcal Q} g_Q(f^{(t)}) d\theta(Q)\right) f^{(t)}
% \end{align*}



% If our domain is curved manifold $\domain$ we can W.L.O.G. (by Takahashi's thm hehe) assume that it's $\domain \in S^{d-1}$, hence, the discrete version of $\Delta_Q$ will require projection $P_i = I - x_i x_i^\top$:
% \begin{align*}
% 	L \in R^{n \times n}: L_{ij}
% 	 & = \exp\left(\frac{-(x_i - x_j)^\top P_i Q^{-1} P_i  (x_i - x_j)^\top}{\epsilon} \right)                                    \\
% 	 & = \exp \left(-\frac{(x_j^\top- x_j^\top x_i x_i^{\top})^\top Q^{-1} (x_j^\top- x_j^\top x_i x_i^{\top})}{\epsilon}\right)
% 	\\
% 	 & = \exp \left(-\frac{(x_j^\top- x_j^\top x_i x_i^{\top})^\top DD^\top (x_j^\top- x_j^\top x_i x_i^{\top})}{\epsilon}\right)
% \end{align*}
% \pagebreak
\newpage

\section{Experiments}



\end{document}